\documentclass[]{../../../../tex/classes/styledArticle}
\usepackage{../../../../tex/packages/hebrewSupport}
\usepackage{../../../../tex/packages/mathShortcuts}
\usepackage{../../../../tex/packages/theoremsSupport}
\usepackage{listings}

\newcommand\din{d_{\mathrm{in}}}
\newcommand\dout{d_{\mathrm{out}}}

\author{שחר פרץ}
\title{\textit{אלגוריתמים 1}}
\date{12 באפריל 2026}
\begin{document}
	\maketitle
	
	\section{הקדמות}
	תרגילים של הפקולטה, יש לנו גישה למודל. 
	
	בקצרה – מה לומדים? הקורס רובו עוסק באלגוריתמים בתורת הגרפים, אך לא כולו. זו כנראה פעם ראשונה שבמקרה אנחנו מתעסקים באיך הגרף מיוצג באמת. זה יותר relaxed מבדידה 1, אין התקטננות על השפה, אך צריך לוודא שהכל נעשה ברור ומסודר, ובאמת מוכיחים את מה שרוצים להוכיח. 
	
	שאלה פשוטה – בהינתן הנחה של 50\% על כל זוג מוצרים, באיזה סדר צריכים להכניס את המוצרים לקופה כדי לקבל את ההנחה המקסימלית? טרוויאלית, סדר מהגבוהה לנמוך יהיה אופטימלי. אבל איך מוכיחים את זה?  קצת טריקי, אבל אפשר להראות שאם מסתכלים על סידור שאיננו האופטימלי, אפשר לבצע החלפה של שני זוגות מוצרים לסידור אחר, שהוא עדיף. 
	
	במילים אחרות, יתקטננו על זה שאפשר לפרמל את ההוכחות, אבל לאף אחד לא באמת איכפת מאוד מהכתיבה. 
	
	נדבר על מציאת עצים פורשים, על משקול, אלגוריתמי זרימה (אפשר לחשוב על זה כמו על צינורות, כאשר הקוטר נקבע ע''י פונ' משקל על הקשתות), אלגוריתמי תכנות לינארי, תכנון דינמי (''שיעור קצת תלוש מהקורס הזה``). בקורס הזה לא מממשים אף מבני נתונים. 
	
	בקורס הזה, יותר מקורסים קודמים, מתעסקים בשאלה גם כמה זכרון אנחנו צורכים, ולא רק סיבוכיות זמן. 
	
	\section{אלגוריתמים}
	מה עושים כאן? מחפשים אלגוריתמים שנכונים לכל \textit{קלט} ונחתשב בזמן הריצה הגרוע ביותר (עבור קלט שמשיג זמן ריצה מקסימלי). כמו בקורסים קודמים, נניח גודל חישוב סטנדרטי וננתח בנוסף לזמן הריצה גם את כמות הזכרון. נעסוק בעיקר באלגוריתמים על גרפים, וכאלה שהם דיטרמינסטיים וסדרתיים (לא נדבר על אלגוריתמים מקביליים, ''ליבה אחת``). שיטות נפוצות: 
	\begin{itemize}
		\item הפרד ומשול. הדוגמה הכי פשוטה, merge sort. פיצול הבעיה לבעיות קטנות, נפתור אותן תחילה ונאחד את הפתרונות. 
		\item אינקרמנטליות: כמו בדוגמה עם המחירים, בונים את הפתרון לאט לאט עד שמקבלים פתרון מלא. 
	\end{itemize}
	
	כמו שאתם זוכרים: 
	\begin{itemize}
		\item גרפים מורכבים \textit{מצמתים/קודקודים} $V$
		\item עם אוסף \textit{צלעות/קשתות} $E$ שמחברים בינהם. 
		
		בגרף מכוון $E$ מורכב מזוגות סדורים, בעוד בגרף לא מכוון $(v, u) \sim (u, v)$ כלומר אפשר לחשוב על זה כמו קבוצות מגודל 2. 
		\item \textit{לולאה} היא קשת מקודקוד לעצמו
		\item \textit{קשתות מקבילות: }אותה הקשת בין אותו זוג הקודקודים
		\item \textit{קשתות אנטי־מקבילות: }(בגרפים מכוונים) קשתות הפוכות בכיוון בין אותם שני הקודקודים. 
		\item \textit{גרף פשוט: }גרף ללא לולאות או קשתות מקבילות. 
		\item \textit{דרגת קודקוד: }כמות הקשתות שחלות בו. בגרפים מכוונים ניתן להגדיר \textit{דרגת כניסה} ו\textit{דרגת יציאה}. 
		\item \textit{תת־גרף: }תת־קבוצה של הקודקודים, $V' \subset V$ ותת קבוצה של הקשתות $E' \subset E$, נותן גרף $(V', E')$. 
		\item \textit{תת־גרף מושרה: }תת גרף בעבורו $E' = (V' \times V') \cap E$. 
		\item \textit{מסלול/מסילה: }סדרת קודקודים וקשתות המחברות בינהן. 
		\item \textit{מעגל: }מסלול שקודקודי ההתחלה והיום שלו זהים. 
		\item \textit{מסלול/מעגל פשוטים: }אין חזרה על קודקודים (פרט לקצוות המעגל). 
		\item \textit{גרף (לא מכוון) קשיר: }ישנו מסלול בין כל שני צמתים. 
		\item \textit{רכיב קשירות: }(בגרף לא מכוון) אלמנט במחלקת השקילות ביחס ליחס שקילות של קשירות בן קודקודים. 
		\item \textit{גרף מכוון קשיר בחוזקה: }גרף מכוון שמכיל מסלול מכוון בין כל שני קודקודים. 
		\item \textit{רכיב קשירות חזקה: }תת גרף מושרה קשיר בחוזקה מקסימלי. 
		\item \textit{מרחק} בין $s$ ל־$v$ הוא מספר הקשתות המינימלי עבורו יש מסלול באורך זה מ־$s$ ל־$v$. נסמן אותו ב־$\dg(s, v)$. אם אין מסלול, נגדיר $\dg(s, v) = \infty$. 
	\end{itemize}
	
	יש לנו מספר טענות שאתם וודאי מכירים: 
	\begin{itemize}
		\item בגרף לא מכוון, סכום הדרגות = פעמיים מספר הקשתות, כלומר $2\sof{E} = \sum_{v \in V}d(v)$
		\item בגרף מכוון, $\sof{E} = \sum_{v \in V}\dout(v) = \sum_{v \in V}\din(v)$
		\item בכל רכיב קשירות חזקה יש מעגל מכוון
		\item בכל גרף לא מכוון, מספר הקודקודים משרדה אי־זוגית הוא זוגי. 
		\item בגרף קשיר ופשוט, מספר הקודקודים חסום לפי $\Omega(\sof{V}) = \sof{V} - 1 \le \sof{E} \le \binom{\sof{V}}{2} = \oc(\sof{V}^{2})$. 
	\end{itemize}
	
	
	\defi{\textit{מעגל אוילר} בגרף קשיר הוא מעגל (לאו דווקא פשוט) שעובר דרך כל קשת בדיוק פעם אחת. }
	האינטואציה אומרת שלכל קודקוד צריך דרגה זוגית – אנחנו עומדים להכנס לקודקוד, ואז לצאת מהקודקוד. 
	\theo{קיים בגרף לא מכוון מעגל אוילר, אמ''מ הגרף קשיר, ודרגת כל הקודקודים היא זוגית. }
	\begin{proof}
		\begin{itemize}
			\item[$\implies$]אם קיים מעגל אוילר, הגרף קשיר, ודרך כל קודקוד נכנסו ויצאנו כל פעם, ולכן דרגתו זוגית. נראה זאת ע''י טיול על המעגל, וספירת הקשתות של כל קודקוד. 
			\item[$\impliedby$]כאן צריך לבנות את המעגל. תיאור מציאת המעגל: 
%			נתחיל מכל קודקוד שבא לנו. נטייל על המעגל, בגלל זוגיות הדרגות נוכל לבצע טיול כזה בלי שהוא יעצר עד שנגיע חזרה לקודקוד שהגענו ממנו. מקשירות, אם לא סיימנו לעבור על כל הגרף, יש קודקוד שלא השתמשנו בו. ''נטייל`` אליו דרך המעגל הראשון בלי להחשיב את זה כחלק ממעגל האוילר, נטייל עליו, ואז נשרשר את שני המעגלים לכדי מהעגל (באינדוקציה אוילר) אחד גדול. מכאן נמשיך עד שנסיים את כל המעגל. 
				נוכיח שקיים מעגל, ע''י מציאתו. נניח שיש לנו גרף לא מכוון, קשיר, וכל הדרגות זוגיות. נתחיל מצומת כלשהו $u \in V$. ניקח שכן שלו $v \in V$, ונוסיף למעגל את הקשת $(u, v)$. מכיוון שהתחלנו עם דרגות זוגיות, כשהגענו ל־$v$ דרגתו ללא הקשת $(u, v)$ נשארה אי־זוגית, ולכן בהכרח נוכל להמשיך ממנו הלאה. נתקע בתהליך רק כשנגיע ל־$u$ בחירה ולא יהיו לו עוד קשתות זמינות. 
				
				כל עוד יש עדיין קשתות בגרף שלא כיסינו, נחפש קודקוד על המעגל שכבר יש לנו עם קשתות זמינות, ונעשה את אותו התהליך ממנו. נחבר (''נשרשר``) את המעגל שקיבלנו למעגל גדול, וחוזר חלילה כל עוד יש עדיין קשתות בגרף. 
		\end{itemize}
	\end{proof}
	אינדוקציה היא בעייתית כאן – תחשבו על מעגל, ברגע שמוציאים קודקוד, הוא לא קשיר יותר. (אפשר לפתור את זה)
	
	תיארנו למעלה אלגוריתם נחמד למציאת מעגל אוילר. אלגוריתם זה עלול להגיע לסיבוכיות של $\omega(\sof{E}^{2})$ בגלל שאנחנו מחפשים כל פעם צומת עם קשתות זמינות. 
	
	ע''מ לחסוך את החיפושים המיותרים אחר הקשתות הזמינות, נחזיק מצביע בודד שמצביע על המעגל, ומחפש קודקוד כזה (שיש לו יותר מ־0 קשתות זמינות). בכל פעם שנתקעים, המצביע לא חוזר אחורה, ולכן כמות הקודקודים/צמתים היא $\sof{E}$. 
	
	\subsection{ייצוג גרפים}
	יש שני ייצוגים סטנדרטיים לגרף – רשימת שכנויות, ומטריצת שכנויות. בעצם, בגרף לעיל הנחנו שאנחנו עובדים ברשימת שכנויות. 
	\begin{itemize}
		\item \textbf{רשימת שכנויות: }רשימה (לא מקושרת) באורך $\sof{V} = n$ שכל איבר בה הוא רשימת הקשתות שיוצאות ממנו (אורך $\sof{E}$ לכל היותר). בגרף לא מכוון, כל קשת תופיע פעמיים. אפשר להניח בלי שום עלות שהרשימות דו־כיווניות, ושיש מצביע בין קשת לקשת ההפוכה שלה. 
		\item \textbf{מטריצת שכנויות: }מטריצה $A \co \sof{V} \times \sof{V}$ כאשר $A_{ij} = I_{(i, j)}$ כאשר $I_x$ האינדיקטור שבודק $x \in E$, אם רוצים לטפל בקשתות מקבילות, במקום $1$ אפשר להכניס את כמות הקשתות. בגרף לא מכוון, מטריצת השכנויות היא סימטרית. 
	\end{itemize}
	הבעיה בשיטה השנייה שבגרפים דלילים, כשאין הרבה קשתות, היא מאוד מאוד לא יעילה. מצד שני, בייצוג מטריציוני נוכל לבדוק ב־$\oc(1)$ האם קיימת קשת. לעומת זאת, למצוא שכנים של קודקוד היא פעולה יקרה (ובאופן כללי, הייצוג דורש $\oc(\sof{V}^{2})$ זכרון, במקום $\oc{\sof{V} + \sof{E}}$). 
	
	\subsection{סריקת גרפים}
	נשאלת השאלה, איך אנחנו יודעים בכלל האם הגרף הוא קשיר? יש לנו שני אלגוריתמים בסיסיים של סריקת גרפים – \en{DFS (death first search)} ו־\en{BFS (breadth first search)}. 
	\subsubsection{BFS}
	נתון גרף לא מכוון בייצוג של מטריצת שכנויות מקודקוד $s$. הלאגו' מבקר בכל הצמתים הנגישים מ־$s$ (ברכיב הקשירות של $s$) בסדר עולה, לפי המרחק מ־$S$ אליהם. האלגוריתם משתמש במבנה נתונים של טור, LIFO. נשמור לכל קודקוד שלושה דברים – color, d, $\pi$. 
	\begin{itemize}
		\item צבע: לבן=לא בוקר, אפור=ראינו ולא סיימנו לטפל בו, ושחור=סיימנו לטפל. 
		\item d=המרחק שגילינו עבור הקודקוד \textit{עד כה}. 
		\item $\pi$=הקודקוד שממנו הגענו לקודקוד הנוכחי. 
	\end{itemize}
	
	להלן השלבים של האלגוריתם: 
	\begin{itemize}
		\item לכל קודקוד $v \in V$, נאפס את הצבע להיות לבן, המרחק לאינסוף וה־$\pi$ ל־NIL. נאפס את הטור $Q$ להיות טור ריק. 
		\item ראשית, נטפל ב־$s$. נגדיר $s.\mathrm{d} \gets 0$ וכן $s.\mathrm{color} \gets \mathrm{gray}$ ולטור נכניס את $s$.
		\item כל עוד $Q$ לא ריק, הקודקודים ב־$Q$ יהיו אפורים, ונבצע את הפעולות הבאות: 
		\begin{itemize}
			\item נגדיר $u = \mathrm{pop}(Q)$ (הראש)
			\item לכל שכן $v$ של $u$: 
			\begin{itemize}
				\item אם $v.\mathrm{color} = \mathrm{white}$, לא ראינו אותו, ולכן: 
				\item נשנה $v.\mathrm{color} \gets \mathrm{gray}$ וכן $v.\mathrm{d} = u.\mathrm{d} + 1$. 
				\item נגדיר $\pi[v] = u$
				\item נדחוף של $u$ ל־$Q$. 
			\end{itemize}
			\item נגדיר את הצבע $u.\mathrm{color} = \mathrm{black}$
		\end{itemize}
		\item נחזיר $d, \pi$. 
	\end{itemize}
	
	בהקשר הזה שחור ואפור זה די אותו הדבר, בגרפים מכוונים יהיה צורך להפעיל if על האפור. התוצאה (אם מסתכלים על ה־$\pi$־ים וחוזרים אחורה) היא עץ פורש ''בשכבות`` שנקבעות לפי $s$. הסיבוכיות: $\oc{\sof{E} + \sof{V}}$. את הלולאה הפנימית עושים $\sum_{v \in V}d(v)$ פעמים לכל היותר, שזה $\oc(\sof{E})$ והאתחול הוא $\oc(\sof{V})$. האיפוס בהתחלה זניח בהנחה שהגרף קשיר. שימו לב שגם ההכנסה וההוצאה מהתור אורכת $\oc(\sof{V})$. 
	
	זה מצב ממש טוב – אלגוריתם שעובד בפחות מ־$\oc(\sof{V} + \sof{E})$ לא עבר על כל הנתונים בגרף. 
	
	$\pi$ מוקדד לנו מסלולים קצרים ביותר מ־$s$ לכל קודקודי הגרף. בעצם, תת־הגרף עם הקשתות $(\pi(v), v)$ לכל $v \in V$ עבורו $\pi(v) \neq \text{NIL}$, הוא עץ מסלולים קצרים ביותר מ־$s$. 
	
	עכשיו נפנה להוכיח שהאלגוריתם עובד, כלומר $v.\mathrm{d} = \dg(v, s)$. 
	
	\theo{טענה: אם $(u, v) \in E$, אז $\dg(s, v) \le \dg(s, u) + 1$ (א''ש המשולש). }\begin{proof}
		פשוט מספיק להוסיף עוד קשת למסלול. 
	\end{proof}
	
	\theo{אם $u$ נמצא קודקוד אחד לפני $v$ במסלול קצר ביותר (מק''ב), מ־$s$ ל־$v$, אז $\delta(s, v) = \delta(s, u) + 1$. }\begin{proof}[הוכחה באמצעות טרוויאלי]
		טרוויאלי
	\end{proof}
	
	\theo{אם $u$ נגיש (יש מסלול) מ־$s$, אזי האלגוריתם יבקר ב־$u$.}\begin{proof}
		באינדוקציה על $k = \delta(s, u)$. בסיס: $k = 0$ כלומר $u = s$ ומתקיים הכל כנדרש. צעד: נניח שנכון עבור $k - 1$ ונוכיח בעבור $k$. קיים מסלול ק''ב מ־$s$ ל־$u$ במרחק $k$. יהא $v$ קודקוד אחד לפני $u$ במסלול. לפי הטענה$\delta(s, v) = \delta(s, u) - 1 = k - 1$ ולכן מההנחה נגיע ל־$v$. מכיוון שיש קשת $(v, u)$, נגיע ל־$u$. 
	\end{proof}
	
	\theo{האלגו' מבקר בצמתים לפי סדר $\mathrm{d}$ עולה. }
	\lem{בכל רגע, ב־$Q$ יש צמתים עם אותו ערך $\mathrm{d}$ או שווי ערי $\mathrm{d}$ עוקבים. }
	''אני דוקטור ואני לא רופא``
	\begin{proof}
		נניח שברגע מסוים ערכי $d$ בתור הם $j, j, j \dots, j + 1, j + 1, \dots$ (הנחת אינדוקציה, הבסיס טרוויאלי), ממויינים. בשלב הבא, נוציא קודקוד מתחילת התור, עבור $\mathrm{d} = j$, ונכניס במקומו קודקודים עם ערך $\mathrm{d} = j + 1$, ולכן השמורה נשמרת. וזה מוכיח את שתי הטענות. 
	\end{proof}
	\theo{לכל $u$ מתקיים ש־$\mathrm{d}[u] \ge \delta(s, u)$ לאורך כל ריצת האלגוריתם. }\begin{proof}
		בהתחלה $\mathrm{d}[u] = \infty$ ולכן מתקיים. עבור קודקוד שלא נגיש מ־$s$, לא נשנה את $\mathrm{d}[u]$. עבור $s$ נאתחל $\mathrm{d}[s] = 0 = \delta(s, s)$ וזה בסדר. 
		
		באינדוקציה על צעדי האלגוריתם, (כאשר $\mathrm{d}[s]$ זה הבסיס) בכל עדכון של $\mathrm{d}[u]$ הוא מקבל את הערך של $\mathrm{d}[w] + 1$ בעבור $w$ הממנו ראינו את $u$, לכן:
		\[ \mathrm{d}[u] = \mathrm{d}[w] + 1 \ge \dg(s, w) + 1 \ge \dg(s, u) \]
	\end{proof}
	
	\theo{בסוף האלגוריתם, $\mathrm{d}[u] = \dg(s, u)$ לכל $u \in U$. }\begin{proof}
		באינדוקציה על $k = \dg(s, u)$. בסיס $k = 0$, $u = s$ מתקיים. צעד ניקח $u$ בעבורו $\dg(s, u) = k$, נניח נכונות עד $k -1$. ניקח מסלול קצר ביותר בין $s$ ל־$u$. יהא $w$ הקודקוד אחד לפני $u$ במסלול. אז: 
		\[ \dg(s, w) \overset{\text{טענה ממקודם}}{=} \dg(s, u) - 1 = k - 1 \overset{\text{induction}}{=} \mathrm{d}[w] \]
		כשהגענו ל־$w$, אך $u$ לבן, אז נקבל: 
		\[ \mathrm{d}[u] = \mathrm{d}[w] + 1 = k - 1 + 1 = \dg(s, u) \]
		מה אם לא? מכיוון שמבקרים בצמתים לפי סדר $\mathrm{d}$ עולה, אך $u$ לא לבן והוא בתור ולכן: 
		\[ k = \dg(s, u) \le \mathrm{d}[u] \le \mathrm{d}[w] + 1 = k \]
		כלומר $\mathrm{d}[u] = k$. 
	\end{proof}
	זה מסיים את הוכחת הנכונות של האלגוריתם. 
	
	אבנה: (בגרף לא מכוון) אם $(u, v) \in E$, אז מתקיים $\mathrm{d}[u] = \mathrm{d}[v] \lor \mathrm{d}[v] = \mathrm{d}[u] + 1 \lor \mathrm{d}[u] = \mathrm{d}[v] + 1$. 
	
	מכאן, שבגרף לא מכוון, קשתות שלא נבחרו לעץ $\pi$ הן באותה השכבה, או בין שתי שכבות עוקבות. בגרף מכוון, מותר לקשתות שלא מקיימות את התנאי הזה לעלות במעלה העץ. 
	
	תזכורת: 
	\defi{גרף $G = (V, E)$ נקרא \textit{דו''צ} אם קיימת חלוקה $V = U \uplus W$ כך שכל הקשתות $(u, w) \in E$ ניתנות לפירוק לכדי $u \in U, w \in W$. }
	
	\theo{גרף לא־מכוון הוא דו''צ אמ''מ אין מעגל אי־זוגי. }\begin{proof}
		\begin{itemize}
			\item[$\implies$]אם הגרף דו''צ כל מעגל מתחיל ומסתיים באותו הצד ולכן באורך אי־זוגי. 
			\item[$\impliedby$]נניח שאין מעגל אי־זוגי. אם הגרף אינו קשיר, נבצע את הפעולות הבאות בעבור כל רכיב קשירות בנפרד (אפשר למצוא אותם גם באמצעות BFS). נריץ BFS על קודקוד $u$ כלשהו. נגדיר את $U$ להיות את כל הקודקודים במרחק זוגי מ־$S$, וב־$W$ להיות כל הקודקודים במרחק אי־זוגי מ־$S$ (או לחלופין המשלים). עבור כל קשת $(u, v)$, מתקיים כמו שראינו ש־$\mathrm{d}[u] = \mathrm{d}[v] \lor \sof{\mathrm{d}[u] - \mathrm{d}[v]} = 1$. נראה שלא ייתכן $\mathrm{d}[u] = \mathrm{d}[v] = k$. במצב זה, קיים בגרף מעגל אי־זוגי $w \to v_k \to v_{k - 1} \to \cdots \to v_1 \to v \to u \to \cdots \to u_k \to w$. באורך $2k + 3$ בסתירה להנחה. 
		\end{itemize}\envendproof
	\end{proof}
	
	\defi{\textit{גרף המסלולים הקצרים ביותר מ־$s$} מכיל את {כל} המסלולים הקצרים ביותר מ־$s$. }
	קבוצה אין־סופית היא קבוצה לא סופית אההההה ווייבס. 
	
	להבדיל מהעץ ש־BFS יורק, גרף המסלולים הקצרים ביותר כולל את כולם, והוא לא בהכרח עץ. 
	
	בהינתן גרף $G = (V, E)$ קשיר ולכן מכוון, וצומת $s$, נרצה למצוא את גרף המקב''ים. 
	
	נוכל לבנות את גרף המקב''ים ע''י הרצת BFS ובחירת כל הקשתות $(u, v)$ כך ש־$\sof{\mathrm{d}[u] - \mathrm{d}[v]} = 1$. 
	
	\textbf{חידה לשבוע הבא: }נאמר ואנחנו מגיעים לכניסה, ואין צורך להקיש על סולמית (כלומר אם הקשנו 12345, בדקרנו גם את 1234). יש 10K קודים. השאלה: כמה לחיצות/הקשות צריכים לעשות כדי לוודא שהדלת נפתחת? 
	
	




	
	
	\ndoc
	
	
	
	
\end{document}